\section{Introduction}
\label{sec:intro}

Parallel token generation can shorten a language model's response time,
but it does not by itself establish a serving advantage. A useful comparison
must account for output length, stopping policy, concurrent requests, and
output quality. Autoregressive (AR) decoding generates tokens sequentially;
block-diffusion decoding refines a canvas of token positions in parallel
before moving to the next canvas. These execution patterns suggest different
latency curves. Their practical consequences require measurement.

We study structured log extraction on a single \gpuModel{} using
vLLM~\cite{vllm2023}. Our primary pair comprises
DiffusionGemma-26B-A4B (\DG)~\cite{diffusiongemma} and
Qwen3-30B-A3B (\ARthree)~\cite{qwen3}. Both are mixture-of-experts
(MoE) models with broadly comparable parameter and weight footprints.
This controls architecture class and approximate scale, but does not match
training data, model quality, tokenizer, or decoder implementation.
The main measurements use the pinned vLLM 0.25.1 stack. Earlier dense
Qwen2.5 and LLaDA experiments are identified separately throughout.

\paragraph{Task and data.}
Each raw log line is mapped to a JSON record with twelve fields. Seven
structural fields describe the timestamp, host, component, severity,
message template, error code, and numeric value; two fields encode an
anomaly label and score; three provide semantic or action information.
The primary corpus is BGL from LOGHUB~\cite{loghub}, containing
4{,}747{,}963 BlueGene/L log lines. The preparation procedure samples
600{,}000 lines by reservoir sampling (seed 42) and splits the sample
temporally into 400{,}000 training and 200{,}000 test records.
Evaluation cells use subsets of this test split, usually $n{=}300$ unless
otherwise stated. HDFS provides a secondary log-corpus check using the
2{,}000-line LOGHUB sample.

Reference records are constructed by the repository's parsing pipeline,
which supports Drain~\cite{drain} and a regex fallback. They are not
independently annotated semantic ground truth. In particular, the code
sets \texttt{is\_anomaly} when the native alert tag is non-normal
\emph{or} severity is ERROR, FATAL, or CRITICAL, and sets
\texttt{anomaly\_score} to the corresponding binary value. Thus these
fields measure agreement with a label-construction rule, not solely
prediction of native alert tags. A FATAL line tagged \texttt{-}, for
example, still receives an anomaly label under this rule.

\paragraph{Metrics and decoding.}
Quality is macro-averaged token-overlap F1 across seven structural and two
anomaly fields. We exclude \texttt{event\_type},
\texttt{affected\_entity}, and \texttt{remediation\_hint}, whose
references are parser-derived rather than independently validated.
Excluding them reduces one source of contamination; it does not make the
remaining labels independent of the preparation pipeline.
The scorer treats two empty field values as a match, so Field-F1 should be
read alongside validity. \emph{Schema validity} is the harness check for
required fields and supported type/enum constraints on the parsed record;
\emph{completeness} measures field presence. Neither establishes semantic
correctness. We use first-object-wins parsing for repeated JSON outputs.

Both models use greedy decoding (temperature 0, top-p 1, repetition
penalty 1). Unless varied explicitly, DG uses $K{=}8$ denoising steps
and canvas size $C{=}256$. Model-specific settings and environment versions
are recorded in the result artifacts. Nominal model sizes are 26B total/4B
active parameters for DG and 30B total/3B active for AR, with approximate
bf16 weight footprints of 52 and 60\,GB; these are rounded model labels,
not measurements of total runtime memory.

We distinguish three quantities: generated output length $L$, the configured
maximum output budget $B$, and the number of decoded canvases
$N_C$. Raising $B$ does not increase $L$ when a model has already emitted
end-of-sequence (EOS). We also distinguish serial mean service cost,
per-event latency percentiles, and inverse throughput under batching.
These quantities answer different deployment questions.

\paragraph{Contributions.}
\begin{itemize}
\item We fit an approximately linear AR cost model and a staircase DG cost
model, specifying the input-length, stopping-policy, and output-range
conditions under which they are useful (\S\ref{sec:laws}).
\item We measure a $2.4\times$ median-latency difference for the primary
pair at batch size one, and report the quality and statistical limitations
of that comparison (\S\ref{sec:wins}).
\item We quantify how batching, latency tails, prompt length, and decoding
effort change the deployment tradeoff (\S\ref{sec:limits}--\ref{sec:survives}).
\item We document nine groups of corrections to our own measurements and
interpretations, including defects that inflated the apparent diffusion
advantage (Appendix~\ref{app:defects}).
\end{itemize}

\begin{table}[ht]
\centering
\caption{Primary comparison. Latency statistics are from the repeated
BGL batch-1 runs (AR: 900 events; DG: 2,000). Throughput multipliers
come from separate batching experiments. DG has greater \emph{relative}
variability but lower \emph{absolute} p99 latency.}
\label{tab:summary}
\begin{tabularx}{\linewidth}{@{}L r r@{}}
\toprule
& \thead{AR-30B} & \thead{DG} \\
\midrule
\tabfont Nominal total / active parameters & 30B / 3B & 26B / 4B \\
\tabfont Approximate bf16 weights & 60\,GB & 52\,GB \\
\tabfont Median latency (ms) & \arMedian & \dgMedian \\
\tabfont p99 latency (ms) & \arPninetyNine & \dgPninetyNine \\
\tabfont Per-event CV (\%) & \arThirtyCV & \dgThirtyCV \\
\tabfont Observed batch-1 to peak multiplier & \arBatchMult$\times$ & up to 1.7$\times$ \\
\bottomrule
\end{tabularx}
\end{table}

\begin{table}[ht]
\centering
\caption{Simplified serial execution models at fixed prompt and decoding
settings. ``Flat'' refers to observed latency within a paid-for canvas,
not zero computation per token.}
\label{fig:schematic}
\begin{tabularx}{\linewidth}{@{}L L L@{}}
\toprule
& \thead{Autoregressive} & \thead{Block diffusion} \\
\midrule
Decoding & Sequential token steps & Parallel positions, sequential canvases \\
Cost approximation & $a_A+b_A L$ & $a_D+s_D(N_C-1)$ \\
Main variable & Actual generated length $L$ & Decoded canvas count $N_C$ \\
\bottomrule
\end{tabularx}
\end{table}
